\scanpage{44}
Revenant à la prop. 1. Soit $C^\bullet$ \inserted{$\in\mathrm{Ob}\,\mathcal B^\bullet$}
un complexe de cochaînes de $\mathcal B$, comment réconstituer le foncteur $k$-lin
\[
 F_{C^\bullet}: \unsure{\mathcal K|\mathcal C|\mathbb K}\longrightarrow\mathcal B
\]
i.e. comment ``calculer'' $F_C(x)$ pour $x\in\mathrm{Ob}\,K$ ?
Il faut qu'on dise comment on ``repère'' $x$,
disons par $\alpha^\bullet(x)\in\mathrm{Ab}_k^\bullet$. Ceci dit, on trouve

\uline{Prop. 3} On a un isom bifonctoriel en $F,x$
(si $F\in\deleted{\unsure{\in|\mathcal B}}\mathrm{Ob}\,\underline{\mathrm{Hom}}_{k\text{-}\mathrm{lin}}(C,\mathcal B)$)
\[
\begin{aligned}
 F(x)&\xrightarrow{\sim}\mathrm{Hom}_k^\bullet(\alpha^\bullet(x),F(C^\bullet))\\
 &\hspace{5.3em}\rotatebox{90}{$=$}\\
 &\hspace{1.3em}\mathrm{Hom}_k^\bullet(\alpha^\bullet(x),C_F^\bullet),
\end{aligned}
\]
[où on définit, pour un $K^\bullet\in\mathrm{Ob}\,\mathrm{Ab}_k^\bullet$
et $C^\bullet\in\mathrm{Ob}\,\mathcal B^\bullet_{\deleted{\unsure{k}}}$,
l'objet $\mathrm{Hom}_k^\bullet(K^\bullet,C^\bullet)$ de $\mathcal B$
(s'il existe) par l'isom fonctoriel en $b\in\mathrm{Ob}\,\mathcal B$
\[
 \mathrm{Hom}(\underset{\underset{\mathrm{Ob}}{\rotatebox{-90}{$\in$}}}{b},
 \mathrm{Hom}_k^\bullet(K^\bullet,C^\bullet))
 \simeq\mathrm{Hom}_{\mathcal B^\bullet}(b\otimes_k K^\bullet,C^\bullet),
\]
si les $K^i$ sont libres de type fini, $b\otimes_k K^\bullet$
est bien défini, donc $b\longmapsto\mathrm{Hom}_{\mathcal B^\bullet}(b\otimes_k K^\bullet,C^\bullet)$
l'est, et la question est si ce foncteur est représentable ; il se trouve qu'il l'est si
$K^\bullet$ est \deleted{\unsure{libre}} $=\in\mathrm{Ob}\,\deleted{\mathrm{Ab}_k}\mathcal K^\bullet$
[ Mais de façon générale, c'est le sous-foncteur de
% slanted
\margin{\unsure{remonter} avant énoncé prop 3}
\[
 \prod_i\underbrace{\mathrm{Hom}_k(K^i,C^i)}_{\in\mathrm{Ob}\,\mathcal B}
\]
noyau de
\[
 \prod_i\mathrm{Hom}_k(K^i,C^i)\mathrel{\underset{\deleted{\unsure{\longrightarrow}}}{\xrightarrow{\delta}}}\prod_i\mathrm{Hom}_k(K^i,C^{i+1})
\]
$(\delta(u^i)=du^i-u^{i+1}d)$ \unsure{...|on} Ça \deleted{est repr.} si
$\mathcal B$ \deleted{\unsure{abélienne}} \inserted{stable par Ker},
si les $\mathrm{Hom}_k(K^i,C^i)$ sont repr (p. ex. \deleted{\unsure{les}} $K^i$ de prés. finie) et les deux
% archive sheet number: 43
